\documentclass[11pt]{article}
\usepackage[margin=1in]{geometry}
\usepackage{amsmath,amssymb,hyperref}
\title{Parallax TRACE Research Signals\\\large Definitions, Construction, and Validation}
\author{Parallax Observatory}
\date{9 November 2026}
\begin{document}\maketitle
\begin{abstract}
This note defines a research-safe construction for signals from FINRA TRACE observations. It separates vendor data, cleaned inputs, descriptive statistics, and candidate forecasts. Nelson--Siegel ticker curves and historical partition yield percentiles are exploratory candidates; neither has been validated as a forecast or investment alpha.
\end{abstract}
\section{Data and eligibility}
For bond $i$ at $t$, let $P_{i,t}$ be clean price per 100 par, $y_{i,t}$ clean yield in decimal units, $D_{i,t}$ modified duration, and $C_{i,t}$ convexity. Preserve instrument/source IDs, trade date, retrieval timestamp, and vintage. The universe $\mathcal U_t$ is an as-of-dated membership set. An observation is eligible only if its policy status and identifiers are valid and it belongs to frozen universe and vintage. Report $q_t=\#\{i\in\mathcal U_t:\mathrm{eligible}\}/\#\mathcal U_t$.
\section{Return primitives}
For successive valid observations, $r^P_{i,t_0,t_1}=\log(P_{i,t_1}/P_{i,t_0})$. Compare elapsed time rather than row count. A yield-space approximation is $\widehat r^Y_{i,t}\approx-D_{i,t^-}\Delta y_{i,t}+\frac12C_{i,t^-}(\Delta y_{i,t})^2$. Yield increases have negative price-return sign. Without duration/convexity, describe yield movement as a rate change, not return.
\section{Cross-sectional residual}
Estimate $\Delta y_{i,t}=\alpha_t+\beta_t^\top x_{i,t^-}+u_{i,t}$ using only contemporaneously available sector, rating, maturity/duration, and liquidity inputs. A candidate dislocation is $z_{i,t}=u_{i,t}/(1.4826\,\mathrm{MAD}(u_{\cdot,t}))$. It becomes a forecast only after out-of-sample estimation of sign and horizon.
\section{Nelson--Siegel ticker curve}
The implemented first candidate groups bonds by the reviewed equity issuer ticker, since rating data are unavailable. For each ticker and as-of date, use each unmatured bond's latest policy-valid TRACE yield with trade date no later than the as-of date. Let $m$ be remaining maturity in years. Fit
\[
y(m)=\beta_0+\beta_1\frac{1-e^{-m/\tau}}{m/\tau}+\beta_2\left(\frac{1-e^{-m/\tau}}{m/\tau}-e^{-m/\tau}\right)
\]
by least squares for $\beta_0,\beta_1,\beta_2$. Select $\tau$ by leave-one-bond-out squared prediction error on 100 log-spaced candidates bounded by the observed maturity range, with a 0.5-year floor and 15-year cap. The default minimum is five bonds and four distinct maturities per ticker. Report fitted yield, observed-minus-fitted yield residual in percentage points, in-sample RMSE, tau-selection leave-one-out RMSE, fitted loadings, and $\tau$. The explorer displays the curve only within the observed maturity range; its value at zero maturity is otherwise unsupported extrapolation when no short bond is present. The leave-one-out score selects $\tau$ and is not an independent performance estimate. A positive residual is a cheap-versus-curve candidate; a negative residual is rich-versus-curve. This cross-sectional fair-yield fit does not itself specify a forecast horizon or establish that residuals mean-revert. Coupon, seniority, optionality, liquidity, rating, and observation staleness can all explain residuals. Walk-forward tests against subsequent bond returns are required before interpreting this as predictive alpha.
\section{Historical partition percentile}
For partition $g$, let $\tilde y_{g,t}$ be the median of eligible latest TRACE yields for its constituent bonds at $t$. The regime signal is the empirical mid-rank percentile of $\tilde y_{g,t}$ against strictly prior partition medians in a trailing lookback window. It is assigned to the bonds in that partition. The test selects whole partition tails; bonds within selected partitions are equally weighted. Future outcomes use raw yield changes rather than cross-sectionally de-meaned changes so that market-wide moves are retained. The current implementation has no benchmark-relative Treasury spread input: yield is only a spread proxy and may reflect rate and maturity-mix changes. Partition membership comes from the current bond-list snapshot, not an as-of universe. This is an exploratory predictor, not validated alpha.
\section{Aggregation and validation}
Group levels may equal-weight eligible bonds. Portfolio signals require explicit weights, turnover, concentration, liquidity, and issuer constraints. Calculate daily return primitives before reporting-period aggregation. Use walk-forward, revision-isolated vintages; report cost-adjusted performance, IC, turnover, capacity, exposures, and regime stability.
\section{Current status}
Parallax stores TRACE price/yield observations and applies an auditable cleaning policy. It displays group levels and rolling level statistics, and calculates ticker-level Nelson--Siegel fitted yields and residuals plus historical partition yield percentiles. These remain exploratory predictors; they are not validated alpha, and the percentile is not a benchmark-relative spread until Treasury curves are joined.

\section{Implementation update: 16 November 2026}
Parallax now persists an immutable maturity-matched Treasury G-spread snapshot, binds it to a TRACE run and Treasury curve vintage, and evaluates both implemented predictors in yield, G-spread, or price space. The walk-forward evaluator uses post-signal changes, reverses yield and spread signs for favourable compression, and retains the log-return sign in price space. Nelson--Siegel outcomes are cross-sectionally demeaned; partition-percentile outcomes are raw. The rebalanced portfolio diagnostic adds an optional one-month Treasury plus G-spread financing proxy and reports price/return coverage.

These are advances in traceability and diagnostics, not validation. The evaluator's default data and funding selections remain ``latest'' artifacts; saved research must pin the cleaned dataset, TRACE run, spread snapshot, curve vintage, bond-list hash, policies, and parameters. Raw-DuckDB and total-return paths that use the compatibility \texttt{observations} table can still be revised by later vendor corrections. Coupon dates are synthetically inferred from maturity. Hence the outputs remain exploratory until pinned research-run manifests, historical attributes, and fixture-based return/curve validation are in place.
\end{document}
